\originalchapter{定态与一维势阱}
\originalsection{定态分离与叠加}
当$V$不含时间, 可寻找分离变量解$\psi(x,t)=u(x)T(t)$. 代入薛定谔方程, 除以$uT$后, 时间变量和空间变量相等, 因而都等于常数$E$. 于是 $\hat Hu=Eu$, $T(t)=\e^{-\ii Et/\hbar}$. 前者是\emph{定态薛定谔方程}, 后者规定这个本征态的时间相位.

定态的密度不随时间改变, 因为$|T|=1$. 多个能量不同的定态叠加一般不是定态. 若 $\psi=\sum_n c_nu_n\e^{-\ii E_nt/\hbar}$, 密度中的交叉项含 $\e^{\ii(E_m-E_n)t/\hbar}$. 动力学的可观测变化正来自相对相位变化. 只求$Hu=Eu$而不展开初始条件, 还没有解决任意初态的时间演化.

\originalsection{无限深势阱}
考虑$0<x<L$中$V=0$, 区间外有无限高墙. 硬墙边界取 $u(0)=u(L)=0$. 在区间内 $u''+k^2u=0$, $k^2=2mE/\hbar^2$. Hamilton算符非负, 因为分部积分给出 $\langle u,Hu\rangle=(\hbar^2/2m)\int|u'|^2\dd x$. $E=0$的线性解在两端为零后只能恒为零, 因此非零本征态有$E>0$.

通解$u=A\sin kx+B\cos kx$. 左边界给$B=0$, 右边界要求$\sin kL=0$, 所以$k=n\pi/L$, $n=1,2,\ldots$. 归一化后,
\[
u_n(x)=\sqrt{\frac2L}\sin\frac{n\pi x}{L},
\qquad E_n=\frac{n^2\pi^2\hbar^2}{2mL^2}.
\]
$n=0$不能作为态: 对应函数是零函数, 无法归一化. 最低能量大于零不是额外的边界约定; 它由局域化和微分方程共同决定. 无限墙处不要求$u'$和外侧零导数连续, 因为这里势能奇异; 若两者都强行设零, 解会消失.

\begin{example}
势阱中初态为$(u_1+u_2)/\sqrt2$. 求其能量测量概率及位置密度.
\end{example}
\begin{solution}
两个能量分别为$E_1,E_2$, 概率各$1/2$, $\langle H\rangle=(E_1+E_2)/2$. 随时间演化后
\[
|\psi(x,t)|^2=\frac{u_1(x)^2+u_2(x)^2}{2}
+u_1(x)u_2(x)\cos\frac{(E_2-E_1)t}{\hbar}.
\]
能量分布不变, 位置分布却振荡; 这两个结论对应不同观测量, 并不矛盾. 正交性使交叉项在整个区间积分为零, 所以总概率仍为1.
\end{solution}

\originalsection{有限深势阱}
取$V=-V_0$在$|x|<a$内, 外部$V=0$, $V_0>0$. 束缚态有$-V_0<E<0$. $E\ge0$时外区振荡、常数或线性解都无法在两端平方可积. 又 $H+V_0$的期望值是非负导数项加非负外区项, $E=-V_0$若要取等只能两项全零, 使解恒为零; $E<-V_0$更不可能. 阱内定义 $k=\sqrt{2m(E+V_0)}/\hbar$, 阱外定义 $\kappa=\sqrt{-2mE}/\hbar$. 外区要取衰减指数, 排除发散指数. 势能有限跳跃没有$\delta$项. 跨边界积分给 $u'(a+\epsilon)-u'(a-\epsilon)=(2m/\hbar^2)\int_{a-\epsilon}^{a+\epsilon}(V-E)u\dd x\to0$, 所以$u'$连续. 若$u$有跳跃,$u''$会出现$\delta'$分布而右端没有对应项, 因而$u$也须连续. 这些条件共同决定能级, 不能只在阱内要求正弦.

反射与对称势的Hamilton算符对易, 所以偶、奇子空间各自保持不变, 可分别寻找本征态. 偶态内部取$A\cos kx$, 右外侧取$B\e^{-\kappa(x-a)}$. 在$x=a$处连续给$B=A\cos ka$, 导数连续给$-\kappa B=-Ak\sin ka$, 所以 $k\tan ka=\kappa$. 奇态内部取$A\sin kx$, 右侧$B\e^{-\kappa(x-a)}$, 左侧须带对应负号. 右边界给 $B=A\sin ka$, $-\kappa B=Ak\cos ka$, 因而$-k\cot ka=\kappa$. 另有$k^2+\kappa^2=2mV_0/\hbar^2$. 这两类超越方程可以图解或数值求根.

令$z=ka$, $z_0=a\sqrt{2mV_0}/\hbar$. 偶态方程是$z\tan z=\sqrt{z_0^2-z^2}$. 在$0<z<\min(z_0,\pi/2)$中, 左边从零递增, 右边递减; 右端或趋零, 或左边趋无穷, 故存在交点. 任意非零深度的一维对称方阱至少有一个束缚态. 三维球阱有不同阈值, 不能把这个结论不加条件搬过去.

外部密度并非零, 而是指数衰减. 经典禁止区表示$E<V$时经典动能无法为正, 不表示量子位置概率必须为零. 束缚态指数尾部没有持续向外的概率流, 所以粒子并不因尾部非零而不断逃走.
